
\medskip
\phantomsection\label{prop:sga3-VIB-8.2}%
\noindent\textbf{Proposition 8.2.---}
\emph{Let \(\mathcal C\) be a category, let \(T\) be a
topology on \(\mathcal C\), let \(G\) be a sheaf of groups on
\(\mathcal C\), and let \(n\) be an integer with \(n\geq0\).
The following conditions are equivalent:}
\begin{enumerate}
\item[\textup{\SGAInlineRefTarget{prop:sga3-VIB-8.2:item:i}(i)}]\phantomsection\label{concl:sga3-VIB-8.2-i}
  \emph{Put \(K_0=G\), and, for \(p\geq1\), put}
  \[
    K_p=\operatorname{Comm}(K_{p-1},K_{p-1})
    \qquad
    \bigl(\textup{respectively, }\
      K_p=\operatorname{Comm}(G,K_{p-1})\bigr).
  \]
  \emph{Then \(K_n\) is the unit presheaf of groups.}

\item[\textup{\SGAInlineRefTarget{prop:sga3-VIB-8.2:item:ii}(ii)}]\phantomsection\label{concl:sga3-VIB-8.2-ii}
  \emph{Put \(K'_0=G\), and, for \(p\geq1\), put}
  \[
    K'_p=\operatorname{Comm}_T(K'_{p-1},K'_{p-1})
    \qquad
    \bigl(\textup{respectively, }\
      K'_p=\operatorname{Comm}_T(G,K'_{p-1})\bigr).
  \]
  \emph{Then \(K'_n\) is the unit sheaf of groups.}

\item[\textup{\SGAInlineRefTarget{prop:sga3-VIB-8.2:item:iii}(iii)}]\phantomsection\label{concl:sga3-VIB-8.2-iii}
  \emph{There is a sequence}
  \[
    G=H_0\supset H_1\supset\cdots\supset H_n
  \]
  \emph{of invariant group subsheaves of \(G\) such that, for
  every \(i\), the quotient sheaf \(H_i/H_{i+1}\) is
  commutative \textup{(respectively, central in
  \(G/H_{i+1}\))}, and \(H_n\) is the unit sheaf.}
\end{enumerate}

It is clear that \(K_n\subset K'_n\); consequently,
\textup{\SGARefLink{prop:sga3-VIB-8.2:item:ii}{(ii)}} implies \textup{\SGARefLink{prop:sga3-VIB-8.2:item:i}{(i)}}. We show that
\textup{\SGARefLink{prop:sga3-VIB-8.2:item:i}{(i)}} implies \textup{\SGARefLink{prop:sga3-VIB-8.2:item:ii}{(ii)}}. One has
\[
  K'_1=\operatorname{Comm}_T(G,G)=K_1^\flat,
\]
and induction using
\hyperref[lem:sga3-VIB-8.1.2]{Lemma 8.1.2} gives
\[
  K'_p=K_p^\flat
\]
for every \(p\). Consequently, if \(K_n\) is the unit
presheaf, then \(K'_n=K_n^\flat\) is the unit sheaf.

Finally, conditions \textup{\SGARefLink{prop:sga3-VIB-8.2:item:ii}{(ii)}} and \textup{\SGARefLink{prop:sga3-VIB-8.2:item:iii}{(iii)}} are
equivalent by
\hyperref[lem:sga3-VIB-8.1.1]{Lemma 8.1.1}.

\medskip
\phantomsection\label{def:sga3-VIB-8.2.1}%
\noindent\textbf{Definition 8.2.1.---}
\emph{When these conditions hold, the sheaf \(G\) is said to
be \textup{solvable of class \(n\)}
\textup{(respectively, nilpotent of class \(n\))}. If there is
an integer \(n\) for which these conditions hold, \(G\) is
said to be \textup{solvable}
\textup{(respectively, nilpotent)}.}

Observe that, by condition \textup{\SGARefLink{concl:sga3-VIB-8.2-i}{(i)}}, this does not depend
on the topology \(T\).

\medskip
\phantomsection\label{prop:sga3-VIB-8.3}%
\noindent\textbf{Proposition 8.3.---}
\emph{Let \(k\) be a field, let \(S\) be a nonempty
\(k\)-scheme, let \(\Omega\) be an algebraically closed
extension of \(k\), and let \(G\) be a smooth \(k\)-group of
finite type. The following conditions are equivalent:}
\begin{enumerate}
\item[\textup{\SGAInlineRefTarget{prop:sga3-VIB-8.3:item:i}(i)}]\phantomsection\label{concl:sga3-VIB-8.3-i}
  \emph{\(G\) is solvable of class \(n\)
  \textup{(respectively, nilpotent of class \(n\))}.}

\item[\textup{(ii)}]\phantomsection\label{concl:sga3-VIB-8.3-ii}
  \emph{\(G\times_kS\) is solvable of class \(n\)
  \textup{(respectively, nilpotent of class \(n\))}.}

\item[\textup{(iii)}]\phantomsection\label{concl:sga3-VIB-8.3-iii}
  \emph{The group \(G(\Omega)\) is solvable of class \(n\)
  \textup{(respectively, nilpotent of class \(n\))}.}

\item[\textup{(iv)}]\phantomsection\label{concl:sga3-VIB-8.3-iv}
  \emph{Put \(K_0=G\), and, for \(p\geq1\), consider the
  \(k\)-groups}
  \[
    K_p=(K_{p-1},K_{p-1})
    \qquad
    \bigl(\textup{respectively, }\ K_p=(G,K_{p-1})\bigr)
  \]
  \emph{\textup{(cf.\ \hyperref[def:sga3-VIB-7.2.2]{7.2.2})}.
  Then \(K_n\) is the unit \(k\)-group.}
\end{enumerate}

The equivalence of conditions \textup{\SGARefLink{prop:sga3-VIB-8.3:item:i}{(i)}} and \textup{\SGARefLink{concl:sga3-VIB-8.3-i}{(ii)}}
follows from
\hyperref[prop:sga3-VIB-8.2]{Proposition 8.2}, since the
formation of the group presheaf of commutators commutes with
base change (IV, \SGARefLink{sga3:iv:proposition:4:1:3}{4.1.3}).

The equivalence of \textup{\SGARefLink{prop:sga3-VIB-8.3:item:i}{(i)}} and \textup{\SGARefLink{concl:sga3-VIB-8.3-iii}{(iv)}} follows
from the fact that, by
\hyperref[prop:sga3-VIB-7.8]{7.8}, the \(k\)-group
\[
  K_p=(K_{p-1},K_{p-1})
  \qquad
  \bigl(\textup{respectively, }\ K_p=(G,K_{p-1})\bigr)
\]
represents the sheaf
\[
  \operatorname{Comm}_T(K_{p-1},K_{p-1})
  \qquad
  \bigl(\textup{respectively, }\
    \operatorname{Comm}_T(G,K_{p-1})\bigr),
\]
where \(T\) is the \textup{(fppf)} or \textup{(fpqc)}
topology.

To prove the equivalence of conditions \textup{\SGARefLink{concl:sga3-VIB-8.3-ii}{(iii)}} and
\textup{\SGARefLink{concl:sga3-VIB-8.3-iv}{(iv)}}, we may assume that \(k=\Omega\); the result then
follows from \hyperref[cor:sga3-VIB-7.10]{7.10}.
